% Exact axis-crossing transfer; all spatial and temporal parameters retained.
\section{An axis-crossing transfer to five dimensions with exact force and energy}
\label{sec:axis-five-dimensional-transfer}

The coordinates in this section are the original Euclidean Cartesian
coordinates \((x_1,x_2,z)\), with \(r=(x_1^2+x_2^2)^{1/2}\) and
\(s=r^2\).  Physical time is \(t\), and viscosity is the prescribed
\(\nu>0\).  In particular \(s=2y_2\) in the preceding meridional
coordinate calculation.  The formulas below extend through the physical
axis.  They give a complete map of smooth equations, their forcing and
their energy; the finite smooth constructions proved here do not assert
a terminal singularity.

\subsection{The spaces and the velocity map, including its inverse}

Let \(I\subset\mathbb R\) be an open time interval.  Take real functions
\(q(s,z,t),h(s,z,t)\) that are smooth on an actual open neighborhood
of \([0,\infty)\times\mathbb R\times I\).  Only their restrictions to
\(s\geq0\) enter any field below.  Assume that there are fixed
\(S,Z>0\) such that these restrictions vanish when \(s>S\) or
\(|z|>Z\), for every \(t\in I\).  They need not vanish at \(s=0\).
This hypothesis fixes the support and guarantees bounded derivatives on
each compact time subinterval.  Define
\begin{equation}\label{eq:aft-velocity}
 \begin{gathered}
 a=-q_z,\qquad c=2q+2s q_s,\\
 u=(x_1a-x_2h,\ x_2a+x_1h,\ c),\qquad s=x_1^2+x_2^2.
 \end{gathered}
\end{equation}
Composition with this polynomial expression for \(s\) proves that \(u\)
is smooth on all of \(\mathbb R^3\times I\), with fixed compact spatial
support.  In the positively oriented cylindrical frame
\(e_r=(x_1/r,x_2/r,0)\), \(e_\phi=(-x_2/r,x_1/r,0)\), \(e_z=(0,0,1)\),
valid for \(r>0\), its components are
\begin{equation}\label{eq:aft-cylindrical}
 u^r=-rq_z,\qquad u^\phi=rh,\qquad u^z=2q+2s q_s.
\end{equation}
The streamfunction and circulation in the earlier convention are exactly
\(\Psi=-s q\) and \(\Gamma=s h\): the formulas
\(u^r=\Psi_z/r\) and \(u^z=-\Psi_r/r\) give
\eqref{eq:aft-cylindrical}.  No cylindrical unit vector is assigned a
value at \(r=0\); the Cartesian formula defines the field there.

Direct differentiation, before using any equation of motion, gives
\begin{equation}\label{eq:aft-divergence-vorticity}
 \begin{gathered}
 \operatorname{div}_3u=2a+2s a_s+c_z=0,\\
 \xi=a_z-2c_s=-(q_{zz}+4s q_{ss}+8q_s),\\
 \omega=\operatorname{curl}_3u
 =(-x_1h_z-x_2\xi,\ -x_2h_z+x_1\xi,\ 2h+2s h_s).
 \end{gathered}
\end{equation}
For example, \(\omega_1=\partial_{x_2}c-\partial_z(x_2a+x_1h)
=-x_2(a_z-2c_s)-x_1h_z\); the other two components follow by the
same differentiation.  Thus \(\omega^\phi=r\xi\) off the axis,
and \(\xi\) is the smooth extension of that quotient.

There is also an inverse, with its domain specified.  Consider fields of
the form \((x_1a-x_2h,x_2a+x_1h,c)\), where \(a,c,h\) have the
smoothness and fixed support just imposed, and
\(2a+2s a_s+c_z=0\).  Put
\begin{equation}\label{eq:aft-inverse}
 q(s,z,t)=\frac12\int_0^1 c(\tau s,z,t)\,d\tau.
\end{equation}
This integral proves smoothness through \(s=0\).  For \(s>0\), it is
equivalent to \(2s q=\int_0^s c(\sigma,z,t)\,d\sigma\), whence
\(c=2q+2s q_s\).  Integration of the divergence identity gives
\(\partial_z\int_0^s c\,d\sigma=-2s a\), so that \(q_z=-a\).
To check support, let \(J(z,t)=\int_0^\infty c(\sigma,z,t)\,d\sigma\).
Then \(J_z=-2[s a]_0^\infty=0\); since \(J=0\) for \(|z|>Z\),
it vanishes for every \(z\).  Therefore \(q=0\) for \(s>S\), as
well as for \(|z|>Z\).  Finally, a difference of two inverse functions
satisfies \(\partial_s(sq)=0\); smoothness at zero forces the resulting
constant \(sq\) to vanish.  This proves the bijection onto the stated
coefficient class, with no unmentioned axis gauge or lost flux.

\subsection{The complete material and diffusion operators}

For smooth functions of \((s,z,t)\), set
\begin{equation}\label{eq:aft-operators}
 \begin{split}
 Lg&=4s g_{ss}+4g_s+g_{zz},\\
 Mg&=4s g_{ss}+8g_s+g_{zz},\\
 Dg&=g_t+2s a g_s+c g_z.
 \end{split}
\end{equation}
The scalar three-dimensional Laplacian is \(L\), whereas
\(\Delta_3(x_jg)=x_jMg\) for \(j=1,2\).  For example
\(\partial_{x_j}g=2x_jg_s\), and summing the differentiated products
gives these two coefficients \(4\) and \(8\).  The full vector advection
is
\begin{equation}\label{eq:aft-advection}
 \begin{split}
 ((u\cdot\nabla)u)_1
 &=x_1(a^2-h^2+2s a a_s+c a_z)
   -x_2(2ah+2s a h_s+c h_z),\\
 ((u\cdot\nabla)u)_2
 &=x_2(a^2-h^2+2s a a_s+c a_z)
   +x_1(2ah+2s a h_s+c h_z),\\
 ((u\cdot\nabla)u)_3&=2s a c_s+c c_z.
 \end{split}
\end{equation}
The terms \(-h^2\) and \(2ah\) have both been retained.

Define the exact residuals
\begin{equation}\label{eq:aft-residuals}
 R_h=Dh-\nu Mh+2ah,\qquad
 E_\xi=D\xi-\nu M\xi-\partial_z(h^2).
\end{equation}
Equivalently, the two scalar equations are
\begin{equation}\label{eq:aft-scalar-equations}
 Dh=\nu Mh+2q_z h+R_h,\qquad
 D\xi=\nu M\xi+\partial_z(h^2)+E_\xi.
\end{equation}
These are exact equations for the original functions, including all
time, convection, stretching and diffusion terms.

\subsection{A force inverse through the axis and the full pressure}

Set
\begin{equation}\label{eq:aft-force}
 F(s,z,t)=\frac12\int_s^\infty E_\xi(\sigma,z,t)\,d\sigma,
 \qquad f=(-x_2R_h,\ x_1R_h,\ F).
\end{equation}
All derivatives of the integrand have fixed compact support.  Differentiation
under the integral is consequently justified on every compact time interval,
and \(F_s=-E_\xi/2\).  More explicitly, near \(s=0\) one can write
\(F=\tfrac12\int_0^\infty E_\xi\,d\sigma
-\tfrac12\int_0^s E_\xi\,d\sigma\) using the given smooth extension;
this proves smoothness on an actual neighborhood of the axis.
The force is zero for \(s>S\) and for \(|z|>Z\).  Its support can fill
the interval between the support of \(E_\xi\) and \(s=0\), which is
still compact in physical space.  In particular, no zero-integral or
mean-zero condition on \(E_\xi\) has been imposed.  Its value at the axis is
\(f^z(0,z,t)=\tfrac12\int_0^\infty E_\xi(\sigma,z,t)\,d\sigma\).
The off-axis cylindrical components and their exact curl relation are
\begin{equation}\label{eq:aft-force-curl}
 f^r=0,\qquad f^\phi=rR_h,\qquad f^z=F,\qquad
 \frac{(\operatorname{curl}f)^\phi}{r}
 =\frac{\partial_z f^r-\partial_r f^z}{r}=-2F_s=E_\xi.
\end{equation}
In Cartesian form the entire curl is
\[
 \operatorname{curl}f=
 (-x_1(R_h)_z-x_2E_\xi,\ -x_2(R_h)_z+x_1E_\xi,
  2R_h+2s(R_h)_s),
\]
so its axis extension is explicit as well.

For completeness, put \(\mathcal F=f+\nu\Delta_3u-u_t-(u\cdot\nabla)u\).
Equations \eqref{eq:aft-operators}--\eqref{eq:aft-residuals} give
\begin{equation}\label{eq:aft-pressure-coefficients}
 \begin{gathered}
 \mathcal F=(x_1A,\ x_2A,\ B),\\
 A=\nu Ma-a_t-a^2-2s a a_s-c a_z+h^2,\\
 B=F+\nu Lc-c_t-2s a c_s-c c_z.
 \end{gathered}
\end{equation}
The azimuthal coefficient vanished because it was exactly
\(R_h+\nu Mh-h_t-2s a h_s-c h_z-2ah=0\).
In particular \(\mathcal F^r/r=A\) is a specified smooth function, even
though the quotient notation itself applies only off the axis.

The closure identity, including its sign, is
\begin{equation}\label{eq:aft-closure}
 A_z-2B_s=-2F_s-E_\xi=0.
\end{equation}
Here is a full algebraic verification.  The time terms are \(-\xi_t\).
The viscosity terms are \(\nu M\xi\), since
\((Lc)_s=M(c_s)\) and \((Ma)_z=M(a_z)\).  The force and swirl terms are
\(-2F_s+\partial_z(h^2)\).  The remaining advection terms are
\[
 N=-\partial_z(a^2+2s a a_s+c a_z)
       +2\partial_s(2s a c_s+c c_z).
\]
Expanding the products groups them as
\[
 N=-2s a\,\partial_s(a_z-2c_s)
     -c\,\partial_z(a_z-2c_s)
     -(a_z-2c_s)(2a+2s a_s+c_z).
\]
The final factor is the divergence already proved to be zero, leaving
\(-2s a\xi_s-c\xi_z\) and proving \eqref{eq:aft-closure}.

Define the pressure, with a fixed time-dependent additive gauge, by
\begin{equation}\label{eq:aft-pressure}
 p(s,z,t)=\int_0^z B(0,\zeta,t)\,d\zeta
              +\frac12\int_0^s A(\sigma,z,t)\,d\sigma.
\end{equation}
Both integrals have smooth integrands through the axis.  They give
\(p_s=A/2\) and, using \(A_z=2B_s\),
\[
 p_z=B(0,z,t)+\frac12\int_0^s A_z(\sigma,z,t)\,d\sigma=B(s,z,t).
\]
Thus the Cartesian gradient is \(\nabla_3p=(x_1A,x_2A,B)=\mathcal F\)
everywhere, including the axis.  We have proved, with the full force and
pressure displayed,
\begin{equation}\label{eq:aft-momentum}
 u_t+(u\cdot\nabla_3)u+\nabla_3p=\nu\Delta_3u+f,
 \qquad \operatorname{div}_3u=0.
\end{equation}
The pressure in this gauge has a constant value on the exterior of the
physical support cylinder, and that constant can be removed explicitly.
Indeed, \(A_z=2B_s\) implies
\(\partial_s\int_{\mathbb R}B(s,z,t)\,dz=0\).  This integral vanishes
for \(s>S\), hence also at \(s=0\); no unequal upper and lower axial
pressure constants arise.  The Cartesian gradient vanishes outside the
closed cylinder \(s\leq S, |z|\leq Z\), whose exterior is connected.
At the exterior point \((s,z)=(S+1,0)\) the value of the pressure is
\[
 C(t)=\frac12\int_0^{S+1}A(\sigma,0,t)\,d\sigma.
\]
Thus \(p-C(t)\) is a smooth pressure supported in that closed cylinder
and gives precisely the same momentum equation.  This conclusion follows
from the constructed force and the closure identity; a pressure-decay
hypothesis was not imposed separately.

\subsection{The five-dimensional operator morphism and its Newton inverse}

Use distinct coordinates \(Y=(X,z)\in\mathbb R^4\times\mathbb R\).
For each scalar coefficient define its lift by
\(\mathcal Jg(X,z,t)=g(|X|^2,z,t)\), and write
\(Q=\mathcal Jq\), \(H=\mathcal Jh\), \(\Xi=\mathcal J\xi\).
This map is injective: every \(s\geq0\) is \(|X|^2\) for some
\(X\).  Its codomain here is exactly the displayed lifts of smooth
coefficient functions, so no regularity claim about a larger invariant
function space is being assumed.  Define
\begin{equation}\label{eq:aft-five-vector}
 V(X,z,t)=(-q_z(|X|^2,z,t)X,\ 2q+2|X|^2q_s).
\end{equation}
Differentiating in the four original \(X\) coordinates proves
\begin{equation}\label{eq:aft-five-intertwining}
 \begin{gathered}
 \Delta_5\mathcal Jg=\mathcal J(Mg),\qquad
 (\partial_t+V\cdot\nabla_5)\mathcal Jg=\mathcal J(Dg),\\
 \Xi=-\Delta_5Q,\qquad
 \operatorname{div}_5 V=4a+2s a_s+c_z=-2q_z=-2Q_z,\\
 V\cdot\nabla_5Q=-2s q_zq_s+cq_z=2q q_z,
 \qquad \operatorname{div}_5(QV)=0.
 \end{gathered}
\end{equation}
Consequently the full transported scalar equations are
\begin{equation}\label{eq:aft-five-pde}
 \begin{split}
 H_t+V\cdot\nabla_5H&=\nu\Delta_5H+2Q_zH+\mathcal J R_h,\\
 \Xi_t+V\cdot\nabla_5\Xi
   &=\nu\Delta_5\Xi+\partial_z(H^2)+\mathcal J E_\xi.
 \end{split}
\end{equation}
This calculation retains the nonzero five-dimensional divergence.
It establishes an exact map of material operators and diffusion, rather
than identifying \(V\) with an incompressible five-dimensional velocity.

The elliptic inverse for these compactly supported coefficients has the
exact constant
\begin{equation}\label{eq:aft-newton}
 Q(Y,t)=\frac{1}{8\pi^2}\int_{\mathbb R^5}
                 \frac{\Xi(Y',t)}{|Y-Y'|^3}\,dY'.
\end{equation}
To prove it, the Gaussian integral in five variables is \(\pi^{5/2}\).
In polar coordinates it is
\(|\mathbb S^4|\int_0^\infty e^{-\rho^2}\rho^4\,d\rho
=|\mathbb S^4|\,3\sqrt\pi/8\), where two integrations by parts
give the last integral.  Therefore \(|\mathbb S^4|=8\pi^2/3\).
For \(\Phi(Y)=|Y|^{-3}/(8\pi^2)\), differentiation shows
\(\Delta_5\Phi=0\) away from zero, and the inward flux for
\(-\Delta_5\) across a sphere is
\(3|\mathbb S^4|/(8\pi^2)=1\).
Green's identity on a punctured ball, followed by shrinking its inner
sphere, gives \(-\Delta_5\Phi=\delta_0\): the term involving
\(\Phi\partial_r\varphi\) has magnitude \(O(\rho)\), whereas the
normal-derivative term converges to \(\varphi(0)\) for each smooth
compactly supported test function \(\varphi\).
The kernel is locally integrable because its radial integral at zero is
\(\int_0^\epsilon\rho\,d\rho\).  Hence, using compact support of
\(Q\), distributional integration by parts gives
\(\Phi*(-\Delta_5Q)=(-\Delta_5\Phi)*Q=Q\), which is exactly
\eqref{eq:aft-newton}.  These equalities can also be tested against a
compact test function and justified by Fubini with the locally integrable
kernel; they introduce no boundary term at infinity.

The compact support in this assertion belongs to \(q\), and therefore
to the particular \(\Xi=-\Delta_5Q\).  It is not a conclusion for an
arbitrary compactly supported right-hand side.  Indeed, if a nonnegative
smooth right-hand side \(G\), supported in \(|Y|\leq R\), has mass
\(m>0\), its Newton potential satisfies for \(|Y|>2R\)
\[
 (\Phi*G)(Y)\geq\frac{m}{8\pi^2(|Y|+R)^3}.
\]
It is consequently neither compactly supported nor rapidly decreasing.
For the actual image of our map, integration by parts additionally gives
\(\int\Xi P\,dY=0\) for every harmonic polynomial \(P\), including
\(\int\Xi\,dY=0\).  These facts describe real restrictions on the
elliptic image; no arbitrary compact \(\Xi\) is silently substituted.

\subsection{The exact measures, energy, dissipation and force work}

The two measures differ.  Angular integration in physical space gives
\(2\pi r\,dr\,dz=\pi\,ds\,dz\).  In \(\mathbb R^5\), integration
over the four-dimensional radial variable gives
\(|\mathbb S^3|r^3dr\,dz=\pi^2s\,ds\,dz\), since
\(|\mathbb S^3|=2\pi^2\).  The latter constant follows, for example,
from \(\pi^2=|\mathbb S^3|\int_0^\infty e^{-r^2}r^3dr
=|\mathbb S^3|/2\).  All integrals in \((s,z)\) below have domain
\([0,\infty)\times\mathbb R\).

The complete physical energy, without a factor \(1/2\) in its definition,
satisfies
\begin{equation}\label{eq:aft-energy}
 \begin{split}
 E_3(t):=\int_{\mathbb R^3}|u|^2\,dx
 &=\pi\int\bigl[s(q_z^2+h^2)+4(q+s q_s)^2\bigr]\,ds\,dz\\
 &=\frac1\pi\int_{\mathbb R^5}(|\nabla_5Q|^2+H^2)\,dY.
 \end{split}
\end{equation}
To see the last equality with all boundary terms, note that
\(|\nabla_5Q|^2=4s q_s^2+q_z^2\), and the difference of the two
meridional integrands is exactly
\[
 4(q+s q_s)^2-4s^2q_s^2=4\partial_s(sq^2).
\]
Its integral is zero at infinity by compact support and at zero by
smoothness and the factor \(s\).  Neither measure was replaced by the
other in this computation.

The complete enstrophy identity is
\begin{equation}\label{eq:aft-enstrophy}
 \int_{\mathbb R^3}|\operatorname{curl}u|^2\,dx
 =\pi\int\bigl[s(\xi^2+h_z^2)+4(h+s h_s)^2\bigr]\,ds\,dz
 =\frac1\pi\int_{\mathbb R^5}(\Xi^2+|\nabla_5H|^2)\,dY.
\end{equation}
The cross terms in the first two Cartesian vorticity components in
\eqref{eq:aft-divergence-vorticity} cancel, and the remaining difference
is \(4\partial_s(sh^2)\), whose two boundary values vanish.  Also
\(\int|\nabla_3u|^2=\int|\operatorname{curl}u|^2\): expansion of the
squared curl gives \(\int|\nabla u|^2-\int\partial_i u_j\partial_j u_i\),
and two integrations by parts identify the latter integral with
\(\int(\operatorname{div}u)^2=0\).  Compact support justifies each
integration.

The force inverse is compatible with these exact measures as follows:
\begin{equation}\label{eq:aft-force-work}
 \int_{\mathbb R^3}u\cdot f\,dx
 =\pi\int s(qE_\xi+hR_h)\,ds\,dz
 =\frac1\pi\int_{\mathbb R^5}
                  (Q\mathcal J E_\xi+H\mathcal J R_h)\,dY.
\end{equation}
Indeed \(u\cdot f=s hR_h+cF\), and
\[
 \int_0^\infty cF\,ds
 =[2s qF]_0^\infty-\int_0^\infty2s qF_s\,ds
 =\int_0^\infty s qE_\xi\,ds.
\]
The boundary at zero vanishes even when \(F(0,z,t)\ne0\), because
\(q,F\) are smooth and the factor is \(s\).

These identities also prove the full energy balance directly from the
five-dimensional equations.  Differentiate \eqref{eq:aft-energy} and
use \(\Xi=-\Delta_5Q\) and its time derivative to obtain
\[
 \tfrac\pi2 E_3'
 =\int_{\mathbb R^5}(Q\Xi_t+HH_t)\,dY.
\]
Substitution of \eqref{eq:aft-five-pde} gives the dissipation
\(-\nu\int(\Xi^2+|\nabla H|^2)\), and the following individual
transport and coupling terms:
\begin{align*}
 -\int Q V\cdot\nabla\Xi&=\int\Xi\operatorname{div}(QV)=0,\\
 \int Q\partial_z(H^2)&=-\int Q_zH^2,\\
 -\int H V\cdot\nabla H&=\tfrac12\int\operatorname{div}V\,H^2
                         =-\int Q_zH^2,\\
 \int H(2Q_zH)&=2\int Q_zH^2.
\end{align*}
Their sum is exactly zero.  The remaining terms are the force work in
\eqref{eq:aft-force-work}.  Division by \(\pi\), followed by
\eqref{eq:aft-enstrophy}, proves
\begin{equation}\label{eq:aft-energy-balance}
 \frac12\frac{d}{dt}\int_{\mathbb R^3}|u|^2\,dx
 +\nu\int_{\mathbb R^3}|\nabla_3u|^2\,dx
 =\int_{\mathbb R^3}u\cdot f\,dx.
\end{equation}
All differentiated integrals are justified by fixed compact spatial
support and bounded derivatives on compact time intervals.

The corresponding local five-dimensional energy flux is also exact.
With \(e_z=(0,0,0,0,1)\), set
\begin{equation}\label{eq:aft-local-flux}
 \begin{split}
 \mathcal B={}&Q\nabla_5Q_t-Q\Xi V
       +\nu(Q\nabla_5\Xi-\Xi\nabla_5Q)\\
       &+QH^2e_z-\tfrac12H^2V+\nu H\nabla_5H.
 \end{split}
\end{equation}
Then, pointwise on all of \(\mathbb R^5\times I\),
\begin{equation}\label{eq:aft-local-balance}
 \partial_t\frac{|\nabla_5Q|^2+H^2}{2}
 +\nu(\Xi^2+|\nabla_5H|^2)
 =Q\mathcal J E_\xi+H\mathcal J R_h+\operatorname{div}_5\mathcal B.
\end{equation}
The product identities proving this formula are
\begin{align*}
 \nabla Q\cdot\nabla Q_t
   &=Q\Xi_t+\operatorname{div}(Q\nabla Q_t),\\
 -QV\cdot\nabla\Xi&=-\operatorname{div}(Q\Xi V),\\
 Q\Delta\Xi&=-\Xi^2+\operatorname{div}(Q\nabla\Xi-\Xi\nabla Q),\\
 Q\partial_z(H^2)&=\operatorname{div}(QH^2e_z)-Q_zH^2,\\
 -HV\cdot\nabla H&=-\operatorname{div}(H^2V/2)-Q_zH^2,\\
 H\Delta H&=-|\nabla H|^2+\operatorname{div}(H\nabla H).
\end{align*}
Here all differential operators in these six lines are five-dimensional.
The first uses \(\Xi_t=-\Delta Q_t\), the second uses
\(\operatorname{div}(QV)=0\), the third uses \(\Delta Q=-\Xi\),
and the fifth uses \(\operatorname{div}V=-2Q_z\).
Adding them with the two scalar equations proves
\eqref{eq:aft-local-balance}, including the cancellation with \(2Q_zH^2\).
Every transverse gradient of a smooth lifted coefficient and every
transverse component of \(V\) is \(O(|X|)\) on compact time intervals.
The flux through an inner tube \(|X|=\epsilon\), over its bounded
axial support, is therefore \(O(\epsilon^4)\) and tends to zero.
This proves explicitly that the local balance has no extra axis flux.

\subsection{Global smooth finite fields and every time-cutoff term}

For explicit global-in-time instances retain arbitrary \(0<T_0<T_1\)
and initial coefficients \(q_0,h_0\) satisfying the hypotheses on an
open interval containing \([0,T_1]\).  Define
\[
 b(\tau)=\begin{cases}e^{-1/\tau},&\tau>0,\\0,&\tau\leq0,\end{cases}
 \qquad
 \eta(t)=\frac{b(T_1-t)}{b(T_1-t)+b(t-T_0)}.
\]
Every derivative of \(b\) for positive \(\tau\) is an exponential
times a polynomial in \(1/\tau\) and tends to zero as \(\tau\downarrow0\).
Thus \(b\), and then \(\eta\), are smooth; the denominator never
vanishes because \(T_0<T_1\).  The values are \(\eta=1\) on
\(t\leq T_0\) and \(\eta=0\) on \(t\geq T_1\).
Take \(q=\eta q_0\), \(h=\eta h_0\) through \(T_1\) and extend
both by zero thereafter.  Flatness at \(T_1\) proves that this extension
is smooth.  Write \(\xi_0=-Mq_0\), let
\(v_0\cdot\nabla_{s,z}=-2s(q_0)_z\partial_s
 +(2q_0+2s(q_0)_s)\partial_z\), and define
\[
 R_0=(h_0)_t+v_0\cdot\nabla h_0-\nu Mh_0-2(q_0)_zh_0,
 \quad
 E_0=(\xi_0)_t+v_0\cdot\nabla\xi_0-\nu M\xi_0-\partial_z(h_0^2).
\]
The complete residuals for the actual cutoff fields are
\begin{equation}\label{eq:aft-time-cutoff}
 \begin{split}
 R_h={}&\eta R_0+\eta'h_0
   +\eta(\eta-1)\bigl(v_0\cdot\nabla h_0-2(q_0)_zh_0\bigr),\\
 E_\xi={}&\eta E_0+\eta'\xi_0
   +\eta(\eta-1)\bigl(v_0\cdot\nabla\xi_0-\partial_z(h_0^2)\bigr).
 \end{split}
\end{equation}
These formulas follow by the product rule: time and diffusion terms
have one factor \(\eta\), and the displayed quadratic terms have
two.  In particular the \(\eta'\) and \(\eta(\eta-1)\) terms are
part of the force \eqref{eq:aft-force}, and have not been discarded.
That integral, together with \eqref{eq:aft-pressure}, now gives a
solution of \eqref{eq:aft-momentum} on the entire time half-line.

All force components and their mixed space-time derivatives are supported
in the fixed cylinder \(s\leq S, |z|\leq Z, 0\leq t\leq T_1\).
Smoothness makes each derivative bounded there.  Consequently, for every
multi-index \(\alpha\), integer \(m\geq0\), and integer \(K\geq0\),
\[
 \sup_{x\in\mathbb R^3,t\geq0}
 (1+|x|+t)^K|\partial_x^\alpha\partial_t^m f(x,t)|<\infty.
\]
The initial velocity has the corresponding Schwartz bounds and is
divergence-free.  To give a completely explicit finite energy estimate,
let \(Q_0,Q_s,Q_z,H_0\) bound respectively
\(|q|,|q_s|,|q_z|,|h|\) on the fixed space-time cylinder.  Then
\begin{equation}\label{eq:aft-uniform-energy-bound}
 \sup_{t\geq0}E_3(t)
 \leq2\pi Z\left\{\frac{S^2}{2}(Q_z^2+H_0^2)
                 +4S(Q_0+S Q_s)^2\right\}<\infty.
\end{equation}
This follows by bounding the two integrands in \eqref{eq:aft-energy}
and integrating \(s\) and \(1\) over \([0,S]\).  The solution itself
is global and smooth, so these instances do not disprove global
regularity.  The established transfer supplies exact axis, material,
elliptic, force, pressure and energy formulas for subsequent constructions
without importing an unproved singular limit.
